\section{Sexta Iteración de la Fase 3: Análisis del problema de las Colisiones}
\label{sec:AnalisisColisiones}

Tenemos al jugador respondiendo a los botones y los gráficos moviéndose por 
''el cristal'', pero hasta ahora los segmentos de jugador y del juego
conviven en pantalla sin interferir unos con otros. 
En esta iteración se analiza un tema muy importante
en el motor del juego, \textbf{el control de colisiones}, que es la 
parte más delicada del motor de estados. 

Como hemos echo hasta ahora, debe ser en el diseñador donde se declaran, 
con la información sobre dónde puede 
haber colisión, y qué pasa si el jugador está allí o no lo está. Las reglas 
estarán definidas de forma que en la iteración siguiente el motor del juego 
use esas reglas \emph{escritas en el cartucho} y solo tenga que
obedecerlas.

Al igual que con las rutas, volveremos al diseñador gráfico para que sea algo
definido por el autor del cartucho (y no constantes fijas de prueba) las que decidan 
una regla de juego. Y como con las rutas, con el mismo principio que
gobierna todo el proyecto: \emph{el cartucho define, el motor ejecuta}.

\subsection{Qué es una colisión en el modelo}

En las máquinas reales no hay cajas que puedan solaparse: el firmware compara
posiciones. Aquí el modelo es de segmentos, y la decisión de diseño es
que la colisión tampoco sea geometría: \textbf{la colisión es un enlace
	declarado}. En un segmento de juego, la autoría marca con qué segmento
del jugador colisiona ese elemento. En tiempo de ejecución, cuando una
ruta activa un segmento que declara colisión, el motor solo tiene que
mirar una cosa: \emph{¿está el jugador visible en el lugar que marca el enlace?}
Si lo está hay \textbf{colisión positiva}; si no está hay \textbf{colisión negativa}. 
Sin cálculo de cajas, sin
solapes, sin umbrales: una comparación de índices. La decisión es del
diseñador del cartucho, y por tanto no tiene ambigüedad posible.

Esta definición encaja con los juegos reales. En el Donkey Kong, el
barril que llega a la fila del jugador mata si el jugador está en ese
lugar; si no lo está, el barril simplemente continúa. En el Parachute,
la aleta del tiburón se traga el bote si está debajo; si no, nada. La
colisión siempre es \emph{elemento en un punto + jugador en ese punto}.

\subsection{Los dos desenlaces de un punto de colisión}

Un punto de colisión es un \textbf{punto de decisión con dos caminos}.
Cuando la ruta activa un segmento del juego, que tiene un enlace al 
segmento de jugador, se genera una colisión, en ese momento con la 
pregunta ``¿está el visible ese segmento?'' decide qué camino sigue el
elemento a partir de ese tick. Para definir correctamente esto, usaremos esta terminología:
\begin{itemize}
	\item \textbf{Colisión Positiva}: El segmento del jugador está visible en el momento de la colisión.
	\item \textbf{Colisión Negativa}: El segmento del jugador NO está visible en el momento de la colisión.
\end{itemize}

En las colisiones, cada uno de estos dos camino puede hacer tres cosas, 
definidas por la autoría del cartucho:

\begin{enumerate}
	\item \textbf{Seguir una cadena}: el elemento deja de seguir
	\texttt{siguiente} y pasa a recorrer el segmento apuntado por el
	enlace del desenlace (la animación del desenlace: las fauces del
	tiburón, el paracaídista rescatado). La cadena se recorre igual
	que una ruta: un segmento por tick, por los enlaces
	\texttt{siguiente} de esa cadena.
	
	\item \textbf{Actualziar un contador}: el desenlace lleva un indicador — vida o muerte — que dice qué contador interno se toca: el de vidas o el de muertes. El contador se actualiza UNA sola vez, en el instante de la decisión: no en cada tick de la animación del desenlace, que puede durar varios segmentos. El indicador pertenece a la decisión, no a la cadena. Los segmentos de tipo 3 y 4 del GWI no SON el contador: son su \textbf{presentación} — qué segmentos se encienden para mostrar el valor del contador es una decisión de presentación (dígitos, marcas, secuencia de encendidos) que ya viaja como parámetro del cartucho, y no cambia en esta iteración..
	
	\item \textbf{Nada}: si el desenlace no declara cadena, el
	elemento simplemente desaparece (igual que muere una ruta al
	llegar al final); si no declara indicador, los contadores no se
	tocan. Los dos son reglas por defecto: lo que no está declarado,
	no pasa.
\end{enumerate}

\subsection{Reutilización de los enlaces del GWI}

El GWI ya viaja con dos enlaces pensados para esto desde la iteración 
en que definimos los enlaces por tipo: \texttt{salvado} y \texttt{muerte}. Esta
iteración les da su semántica definitiva, sin crear campos nuevos ni
romper los cartuchos existentes:

\begin{itemize}
	\item \texttt{salvado} es el desenlace \textbf{con jugador}: qué
	cadena sigue el elemento cuando el jugador SÍ está en el punto de
	colisión.
	
	\item \texttt{muerte} es el desenlace \textbf{sin jugador}: qué
	cadena sigue cuando el jugador NO está.
\end{itemize}

Los nombres son los del dominio y vienen de la lógica de las máquinas
reales — ``si el bote está, el tiburón se lo traga (muerte); si no, la
presa se salva (salvado)'' — pero conviene dejar escrito que el motor
no interpreta los nombres: el motor solo sigue el enlace que le toca
según esté o no visible el jugador. La \emph{vida} o \emph{muerte} del
desenlace la dice el indicador, no el nombre del enlace.

¿Y si un desenlace es de vida o de muerte sin cadena (sin animación)?
Las reglas por defecto lo cubren: el elemento desaparece y el
indicador, si lo hay, toca su contador. Un paracaidista que colisiona 
en positivo desaparece y aumenta el contador de vidas, la colisión en 
negativo lanza la ruta de muerte, que al final tiene su propio contador.

\subsection{El indicador de desenlace}

Junto al enlace, cada desenlace lleva un indicador opcional con tres
valores posibles:

\begin{itemize}
	\item \textbf{vida}: la colisión aumenta el contador de vidas (presentado por los segmentos de tipo 3).
	\item \textbf{muerte}: la colisión aumenta el contador de muertes (presentado por los segmentos de tipo 4).
	\item \textbf{nada} (o ausente): los contadores no se tocan; el
	desenlace es solo animación.
\end{itemize}

El indicador viaja en el GWI junto a su enlace
(\texttt{salvadoIndicador} / \texttt{muerteIndicador}), porque es
propiedad del desenlace: \texttt{salvado} puede ser de vida y
\texttt{muerte} de muerte, o al revés, o ambos de nada. La autoría
decide. Un punto de colisión \emph{con} los dos enlaces definidos y
\emph{sin} ninguno de los dos indicadores es válido: será puro teatro
(el elemento cambia de cadena según esté el jugador, pero nadie muere).

\subsection{Qué va a tocar este requisito}

Se hará en dos partes, en una iteración se añaden los cambios necesarios 
en el diseñador gráfico, en la siguiente se aborda el cambio del motor
de estados. El motor gráfico no cambia. Concretamente:

\begin{itemize}
	\item \textbf{Panel}: los segmentos de tipo juego ganan la marca
	de colisión — seleccionar el segmento del jugador con el que
	colisiona — y, para cada uno de los dos desenlaces, su cadena
	(el enlace ya existente) y su indicador (vida/muerte/nada).
	
	\item \textbf{Mapa de objetos}: los segmentos que declaran
	colisión se señalan con un color propio (\texttt{teal}), y el
	trazo de la relación colisión—jugador se pinta también, para ver
	de un vistazo dónde están los puntos de colisión del cartucho.
	
	\item \textbf{Verificación}: los mismos criterios que ya se
	aplican a los enlaces del jugador — el inverso calculado se
	guarda en el GWI (el segmento jugador sabe qué puntos de
	colisión dependen de él), avisos de conflictos (un jugador sin
	ningún punto de colisión asociado, un desenlace con cadena pero
	sin colisión declarada), y la exigencia de que los segmentos
	que deben colisionar (los de la última línea de las rutas del
	cartucho de prueba) tengan su colisión declarada antes de
	poder guardar.
\end{itemize}

\subsection{Decisiones menores}

\begin{itemize}
	\item Todo segmento que declara colisión debe tener
	ambos desenlaces definidos, si uno no lo está se usa la
	regla por defecto: \textbf{desaparece sin tocar contadores}.
	
	\item Un segmento de jugador no puede ser punto de
	colisión de varios elementos, se deja como posible
	ampliación que se permita, aunque no se ve este caso en los
	juegos que se han consultado.
	
	Pero no hay que confundir la cadena de desenlace con la ruta de juego: las aletas del tiburón bajo las barcas en el Parachute no son segmentos de ruta — forman parte del desenlace de muerte o de una animación independiente — y por eso no declaran colisión.
	
	
	\item La verificación avisará de una colisión declarada
	con un segmento del jugador \emph{inalcanzable} por los enlaces
	de dirección del jugador, bloqueado la creación del cartucho hasta
	su resolución. No permitimos datos incoherentes.
\end{itemize}

\clearpage